\documentclass[11pt]{article}
\usepackage[T1]{fontenc}
\usepackage{lmodern}
\usepackage{amsmath,amssymb,amsthm,mathtools,booktabs,array}
\usepackage[margin=1in]{geometry}
\usepackage[colorlinks=true,urlcolor=blue,citecolor=blue]{hyperref}
\newtheorem{theorem}{Theorem}
\newtheorem{lemma}[theorem]{Lemma}
\newtheorem{corollary}[theorem]{Corollary}
\newtheorem{proposition}[theorem]{Proposition}
\theoremstyle{remark}
\newtheorem{remark}[theorem]{Remark}
\newcommand{\rem}{\operatorname{rem}}
\newcommand{\calls}{\operatorname{calls}}
\newcommand{\eq}{\operatorname{eq}}
\newcommand{\coprime}{\perp\!\!\!\perp}
% Lean identifiers, breakable after dots and underscores.
\newcommand{\lean}[1]{\nolinkurl{#1}}
\makeatletter\g@addto@macro\UrlBreaks{\do\_}\makeatother
\emergencystretch=2em
\title{The Euclidean algorithm is not optimal among uniform processes}
\author{Jason Davies\thanks{\href{mailto:jason@jasondavies.com}{\nolinkurl{jason@jasondavies.com}}}}
\date{October 23, 2026}
\begin{document}
\maketitle
\begin{abstract}
Reusing a stored Euclidean prefix, we derive $\gcd(a,b)$ from $a,b$ with
$O(\log a/\log\log a)$ applications of the remainder function alone,
uniformly for $a>b\geq1$; Euclid's algorithm needs order $\log a$ in the
worst case. These derivations bound the intrinsic complexity
of gcd and coprimeness by the same order, refuting Moschovakis's conjectured
logarithmic lower bounds for uniform processes; by his best-uniform-process
theorem, a single uniform process attains this bound.
The Main Conjecture for algorithms remains open. We also prove an
$\Omega(\log\log a)$ lower bound for almost all coprime pairs and give a
family of finite nondeterministic gcd programs with arbitrarily small
logarithmic cost coefficients. The bounds and the refutations are
formalised in Lean~4; the best-uniform-process theorem is cited.
\end{abstract}

\noindent\textbf{AI transparency statement.} This paper was produced entirely
by artificial intelligence, including the identification of the main result,
its formalisation in Lean, and the writing of the manuscript. The human author
initiated the broader research and directed the AI systems, but did not write
the text. Details are given at the end of the paper.

\section{Introduction}\label{sec:intro}
Van den Dries and Moschovakis asked whether the Euclidean algorithm is optimal
among its peers, the algorithms that compute the gcd from the remainder
function \cite[p.~390]{vddm2004}. Moschovakis's Main Conjecture asserts that for every
algorithm $\alpha$ from $\rem$ and the zero test $\eq_0$ computing
$\gcd(x,y)$ there is an $r>0$ such that $\alpha$ makes more than $r\log x$
remainder calls on infinitely many pairs $x>y\geq1$
\cite[p.~2]{moschovakis2019}. We refute its formulations for uniform
processes, the book's more general model (Section~\ref{sec:certificates});
for algorithms it remains open.
Throughout, ``the book'' means Moschovakis's monograph
\cite{moschovakis2019}, and logarithms are to base $2$.

\paragraph{Derivations.}
A \emph{derivation} from $a,b$ is a finite list that starts with $a,b$ and in
which each further entry is $\rem(x,y)$ for earlier entries $x,y$ with $y>0$.
Each further entry is one \emph{instruction}, that is, one remainder call.
Entries already in the list are \emph{available} and may be reused; no
constants and no other operations are supplied. Let $G(a,b)$ be the least
number of instructions in a derivation containing $\gcd(a,b)$, and let
$D(a,b)=G(a,b)$ for coprime inputs. Thus $G(a,0)=0$ and $D(a,1)=0$.
Arithmetic used in the proofs, such as subtraction, sorting or comparison,
only selects instructions; it is never itself an instruction. Thus $G$
measures the length of a derivation, not the work of finding one.

\begin{theorem}\label{thm:main}
There is a constant $C$ such that for all integers $a\geq4$ and $0\leq b\leq a$,
\[
G(a,b)\leq C\,\frac{\log a}{\log\log a}.
\]
In particular this bound holds for $D(a,b)$ when $\gcd(a,b)=1$.
\end{theorem}

\paragraph{Idea of the proof.}
For sufficiently large $n$, let $a>b$ be coprime with $a\leq N=2^n$.
Store a Euclidean prefix
$r_0=a,r_1=b,\ldots,r_h$, where $h=\lceil\sqrt n\rceil$.
If the prefix reaches a positive remainder at most $N^{3/4}$, take one
further remainder and recurse on the resulting smaller pair.
Otherwise its consecutive pairs stay distinct modulo every $m>N^{1/3}$, so
the prefix occupies at least $\sqrt h$ residue classes modulo $m$
(Lemma~\ref{lem:prefix}). A close pair of residues lets three remainder
calls produce a positive value below $m/(\sqrt h-1)$
(Lemma~\ref{lem:gap}). Reusing the same prefix, we repeatedly shrink the
modulus from $r_h$ to at most $N^{1/3}$ in $O(n/\log n)$ calls.
Two recursive gcd derivations on inputs at most $N^{1/3}$ finish.
If $T(n)$ is the worst derivation cost for inputs at most $2^n$, the
resulting recurrence \eqref{eq:recurrence} gives $T(n)=O(n/\log n)$.
Tuning the prefix length sharpens the upper coefficient to $6+o(1)$
(Corollary~\ref{cor:six}).

\paragraph{Consequences and further results.}
Uniform processes are the book's abstract model of computation from given
primitives; every deterministic or nondeterministic recursive program
induces one \cite[\S4A, Proposition~4A.1]{moschovakis2019}.
Section~\ref{sec:certificates} converts the derivations into small
certificates. This bounds the intrinsic complexity of gcd and coprimeness,
refutes Conjecture~3.4 of \cite{moschovakis2012} and version (4F-5) of the
Main Conjecture, and limits what the homomorphism method can prove. The
book's best-uniform-process theorem then gives a single process with the
same bound on every input. The upper bound also settles an open problem
of van den Dries and Moschovakis on the closure depth of the gcd
(Corollary~\ref{cor:closure-depth}).

We also prove that $D(a,b)\geq(1-o(1))\log\log N$ for all but $o(N^2)$
coprime pairs $2\leq b<a\leq N$ (Proposition~\ref{prop:loglog}).
Finally, fixing the prefix length gives finite nondeterministic programs
with arbitrarily small logarithmic coefficients, uniformly over all inputs
(Theorem~\ref{thm:programs}). Their program size grows as the requested
coefficient decreases. They need no recursion: register programs with a
fixed number of registers suffice (Theorem~\ref{thm:registers}).

\paragraph{Related work.}
The known intrinsic lower bounds give $\Omega(\log\log a)$ on special
families of coprime pairs \cite{vddm2004,vddm2009},
\cite[Theorem~6C.5]{moschovakis2019}, also for richer primitives such as
addition and integer quotient, and for depth. Pratt's unpublished nondeterministic algorithm
\cite[Theorem~2E.2 and Problems~x2E.17--x2E.19]{moschovakis2019}
matches this order on successive Fibonacci numbers
\cite[Corollary~6C.7]{moschovakis2019}.
For remainder-only derivations, our almost-all bound gives the same order
on a set of full density among coprime pairs.

In a different model, Bshouty obtained the optimal bound
$\Theta(\log N/\log\log N)$ for gcd on inputs in $[0,N]$ using a random
access machine with linear operations, shifts, comparisons and indirect
addressing \cite[Theorems~i and~ii, p.~3]{bshouty1989}. His construction
tabulates blocks of the binary gcd algorithm \cite{stein1967},
\cite[\S4.5.2]{knuth1997}, computing table entries on demand and counting
their construction \cite[\S4, Algorithm~D and Lemma~5, pp.~12--14]{bshouty1989}.
Here the saving comes from
reusing a Euclidean prefix with remainder as the only instruction.

\section{Normalisation}\label{sec:normalization}
Write $g=\gcd(a,b)>0$. For positive $y$,
\[
\rem(gx,gy)=g\rem(x,y).
\]
Every generated value is divisible by $g$. Consequently scaling a derivation,
or dividing all its values by $g$, preserves its length and the validity of
every instruction. Hence
\begin{equation}\label{eq:normalize}
G(a,b)=D(a/g,b/g).
\end{equation}
Define
\[
T(n)=\max_{0\leq b\leq a\leq 2^n,\ a>0}G(a,b).
\]
Euclid's remainders form a derivation containing $\gcd(a,b)$, so $G$ and
$T$ are finite. By \eqref{eq:normalize} it
suffices to bound $D$; the cases $a=b$, $b=0$ and $b=1$ need no instructions.

\section{Three-call contractions}\label{sec:contractions}
\begin{lemma}\label{lem:gap}
If $m>0$ is available and a finite collection $S$ of available values occupies
$t\geq2$ residue classes modulo $m$, at most three instructions generate $w$
with
\[
0<w<\frac{m}{t-1}.
\]
\end{lemma}
\begin{proof}
Sort the distinct residues. Some adjacent pair $u<v$ has gap
$\delta=v-u<m/(t-1)$, because the sum of the $t-1$ adjacent gaps is less than
$m$. Each selected residue is $\rem(x,m)$ for some $x\in S$, so it costs one
instruction. If $u=0$, use $v$. If $0<u\leq\delta$, use $u$. Otherwise
$0<\delta<u$ and $v=u+\delta$, so $\rem(v,u)=\delta$. This uses at most three
instructions.
\end{proof}

\begin{lemma}[Euclidean-prefix lemma]\label{lem:prefix}
Let $r_0=a>r_1=b>\cdots>r_h>0$ be a Euclidean prefix of a coprime pair, with
$r_{i-1}=q_i r_i+r_{i+1}$ and $q_i\geq1$ for $1\leq i<h$. For every integer
$m>2a/r_h$, the
$h$ pairs
\[
(r_i\bmod m,r_{i+1}\bmod m),\qquad 0\leq i<h,
\]
are distinct. The prefix therefore occupies $t$ residue classes with $t^2\geq h$.
\end{lemma}
\begin{proof}
Suppose the pairs at $i<j$ coincide. Set
\[
M=\prod_{\ell=i+1}^{j}\begin{pmatrix}q_\ell&1\\1&0\end{pmatrix}
 =\begin{pmatrix}A&B\\C&E\end{pmatrix},\qquad
\binom{r_i}{r_{i+1}}=M\binom{r_j}{r_{j+1}}.
\]
Right multiplication by
$\bigl(\begin{smallmatrix}q&1\\1&0\end{smallmatrix}\bigr)$ sends
$(A,B,C,E)$ to $(qA+B,\,A,\,qC+E,\,C)$, starting from $(q,1,1,0)$ for a single
factor. By induction, all entries are nonnegative, $A\geq B\geq E$,
$A\geq C\geq1$, $B\geq1$, and $\det M=(-1)^{j-i}$.
Also $Ar_j\leq r_i\leq a$ and $r_j>r_h$ (as $j<h$), whence $A<a/r_h$.
Write
\[
\Delta=\det(M-I)=1+\det M-A-E.
\]
For an odd number of factors, $\Delta=-A-E<0$. By the same recursion, a
product of at least two factors has $A\geq2$ and $E\geq1$; thus for an even
positive number of factors,
$\Delta=2-A-E<0$. In either case,
\[
0<|\Delta|\leq A+E\leq2A\leq2a/r_h<m.
\]
The presumed congruence gives $(M-I)(r_j,r_{j+1})^{\mathsf T}\equiv0\pmod m$.
Multiplication by the adjugate shows that $m$ divides both $\Delta r_j$ and
$\Delta r_{j+1}$. Consecutive remainders are coprime, so B\'ezout's identity
implies $m\mid\Delta$, a contradiction. Finally, $t$ individual residue classes
give at most $t^2$ ordered pairs.
\end{proof}
Since Lemma~\ref{lem:prefix} holds for every $m>2a/r_h$, one stored prefix
serves every later modulus, although these moduli are not known in advance.
When $\sqrt h>1$, Lemma~\ref{lem:gap} applied to the prefix gives three
instructions that generate
\begin{equation}\label{eq:contract}
0<w<\frac{m}{\sqrt h-1}.
\end{equation}
We call these three instructions a \emph{contraction} of $m$, and $w$ the new
modulus.

For example, from $44,17$ one instruction stores the prefix $44,17,10$.
Modulo $m=10>2\cdot44/10$ its residues are $4,7,0$, and the close pair
$4<7$ gives the contraction
\[
\rem(44,10)=4,\qquad \rem(17,10)=7,\qquad \rem(7,4)=3,
\]
with new modulus $3<10/(3-1)$, as in Lemma~\ref{lem:gap}. One more
instruction, $\rem(10,3)=1$, finishes. Euclid's algorithm also takes five
instructions here; the saving appears only for large inputs.

\section{Construction and recurrence}\label{sec:construction}
Fix an integer $n>16$ and put $N=2^n$ and $h=\lceil\sqrt n\rceil>4$.
After normalisation, assume $a>b>1$ are coprime and $a\leq N$.
We use $N^{3/4}$ as the \emph{early threshold} and $N^{1/3}$ as the
\emph{target modulus}. Any thresholds $N^\alpha,N^\beta$ with $\alpha<1$,
$\alpha+\beta>1$ and $2\beta<1$ would do: these make the early recursive
problem smaller, let the prefix lemma apply above the target, and make the
two final problems together smaller. We take $\alpha=3/4$ and $\beta=1/3$.

\paragraph{Phase (i): the prefix.}
Starting from $r_1=b$, generate the Euclidean remainders $r_2,r_3,\ldots$,
one instruction each, until the most recent value is at most $N^{3/4}$ or
$r_h$ is available. This costs at most $h$ instructions. Since the inputs
are coprime, Euclid's algorithm reaches $1$ before $0$, so the stopping
value is positive.

If the stopping value $c$ is at most $N^{3/4}$, generate the next
remainder $s<c$ with one more instruction. The pair $(c,s)$ is coprime, with
both entries at most $2^{\lceil3n/4\rceil}$. This gives the bound
\begin{equation}\label{eq:early}
h+1+T(\lceil3n/4\rceil).
\end{equation}

\paragraph{Phase (ii): contractions.}
Otherwise $r_h>N^{3/4}$, so $2a/r_h<2N^{1/4}<N^{1/3}$ (as $n>12$), and
Lemma~\ref{lem:prefix} applies to every available
$m>N^{1/3}$. Starting at $m=r_h$, repeatedly replace $m$ by the value $w$ of
\eqref{eq:contract} until $0<m\leq N^{1/3}$. Set
\[
\kappa=\sqrt h-1>1,\qquad
J=\left\lceil\frac{2n/3}{\log\kappa}\right\rceil.
\]
There are at most $J$ contractions: if $J$ have occurred then
$m<N/\kappa^{J}\leq N^{1/3}$. Since
$\log\kappa=\tfrac14\log n+O(1)$, their cost is $3J=O(n/\log n)$.
This is the source of the saving: the prefix costs only $O(\sqrt n)$
instructions, and each contraction reduces $\log m$ by at least
$\log\kappa=\tfrac14\log n+O(1)$.

\paragraph{Phase (iii): two recursive derivations.}
The small $m$ need not be coprime to either input. Generate
$u=\rem(a,m)$ and $v=\rem(b,m)$, then append a derivation of $d=\gcd(m,u)$
from $m,u$ and one of $e=\gcd(d,v)$ from $d,v$. Since
$m\leq2^{\lceil n/3\rceil}$ and $G$ is symmetric, each costs at most
$T(\lceil n/3\rceil)$. Finally
\[
e=\gcd(\gcd(m,\rem(a,m)),\rem(b,m))
 =\gcd(m,\gcd(a,b))=1.
\]
The cost in this case is at most
\begin{equation}\label{eq:late}
h+3J+2+2T(\lceil n/3\rceil).
\end{equation}

\paragraph{The recurrence.}
The two case estimates \eqref{eq:early} and \eqref{eq:late} and normalisation
imply
\begin{equation}\label{eq:recurrence}
T(n)\leq\max\{T(\lceil3n/4\rceil)+O(\sqrt n),\;
2T(\lceil n/3\rceil)+O(n/\log n)\}.
\end{equation}
All implied constants are independent of the inputs.

\begin{proof}[Proof of Theorem~\ref{thm:main}]
For $n\geq2$ let $f(n)=n/\log n$. The ratios
\[
\frac{f(\lceil3n/4\rceil)}{f(n)}\longrightarrow\frac34,
\qquad
\frac{2f(\lceil n/3\rceil)}{f(n)}\longrightarrow\frac23
\]
show that, beyond a fixed threshold, $f(\lceil3n/4\rceil)\leq\tfrac78f(n)$ and
$2f(\lceil n/3\rceil)\leq\tfrac56f(n)$. Both overheads in
\eqref{eq:recurrence} are at most $c_0f(n)$ for a fixed $c_0$. Choose the
induction threshold above $16$, where the recurrence holds, and large enough
that both recursive indices are at least $2$ and strictly less than $n$. Choose
$C\geq8c_0$, increased to cover the finitely many values below the threshold.
Strong induction then gives $T(n)\leq Cf(n)$: the two bounds are
$(7C/8+c_0)f(n)$ and $(5C/6+c_0)f(n)$, both at most $Cf(n)$.
Finally take $n=\lceil\log a\rceil$. For $a\geq4$, $n\leq2\log a$ and
$\log n\geq\log\log a\geq1$, so $G(a,b)\leq T(n)\leq2C\log a/\log\log a$.
\end{proof}

\subsection{The leading coefficient}
The same construction allows a more precise bound without changing either
arithmetic lemma.

\begin{corollary}\label{cor:six}
For every $\varepsilon>0$, there is an $a_0$ such that for every
$a\geq a_0$ and $0\leq b\leq a$,
\[
G(a,b)\leq(6+\varepsilon)\frac{\log a}{\log\log a}.
\]
Equivalently, the uniform upper bound has coefficient $6+o(1)$.
\end{corollary}
\begin{proof}
Fix $0<c<1/2$ and $0<\beta<1/2$, and put $\alpha=1-\beta/2$.
For large $n$ put
\[
p=\lfloor c\log n\rfloor,\qquad k=2^p,\qquad h=k^2+1,
\qquad J=\lceil n/p\rceil.
\]
Use $N^\alpha$ as the early threshold and $N^\beta$ as the target modulus,
where $N=2^n$. In the late case $r_h>N^\alpha$, so
$2a/r_h<2N^{1-\alpha}=2N^{\beta/2}<N^\beta$ for sufficiently large $n$, and
the prefix lemma applies to every modulus above $N^\beta$. Since $k^2<h$, the prefix occupies at least $k+1$ residue classes,
so each contraction reduces the modulus by a factor greater than $k$.
Moreover $k^J\geq N$, so $J$ contractions suffice.
Writing $f(n)=n/\log n$, we have
\[
h=O(n^{2c})=o(f(n)),\qquad
3J=\left(\frac3c+o(1)\right)f(n).
\]
The two branches therefore give
\[
T(n)\leq\max\left\{
 T(\lceil\alpha n\rceil)+o(f(n)),\quad
 2T(\lceil\beta n\rceil)+\left(\frac3c+o(1)\right)f(n)
\right\}.
\]
Theorem~\ref{thm:main} implies that
$\Lambda=\limsup_{n\to\infty}T(n)/f(n)$ is finite and nonnegative.
For each fixed $q>0$, the ratio $f(\lceil qn\rceil)/f(n)$ tends to $q$.
Taking upper limits in the recurrence yields
\[
\Lambda\leq\max\{\alpha \Lambda,\ 2\beta \Lambda+3/c\}.
\]
Since $\alpha<1$ and $\Lambda\geq0$, this implies
$(1-2\beta)\Lambda\leq3/c$. First let $\beta\downarrow0$ and then $c\uparrow1/2$ to obtain
$\Lambda\leq6$. Finally, $n=\lceil\log a\rceil$ satisfies
$f(n)\sim\log a/\log\log a$, proving the assertion. All upper limits concern
the worst-case function $T$, so the resulting threshold is independent of $b$.
\end{proof}
We do not claim that $6$ is optimal.

\subsection{Closure depth}
Van den Dries and Moschovakis \cite{vddm2009} measure the parallel depth of the gcd. Put
$S_0=\{0,1,a,b\}$, let $S_{k+1}$ adjoin to $S_k$ all values $x+y$, $x-y$,
$\operatorname{iq}(x,y)$ (integer quotient) and $\rem(x,y)$ with
$x,y\in S_k$, and let $g(a,b)$
be the least $k$ with $\gcd(a,b)\in S_k$ \cite[p.~2:6]{vddm2009}. They ask
whether there are positive integers $\alpha_n\to\infty$ and $\beta_n$ and a
$C>0$ with $g(\alpha_n,\beta_n)>C\log\alpha_n$ for all $n$
\cite[Open Problem, p.~2:11]{vddm2009}.

\begin{corollary}\label{cor:closure-depth}
For integers $\alpha\geq4$ and $\beta\geq1$,
\[
g(\alpha,\beta)\leq1+G(\alpha,\rem(\beta,\alpha))
 =O\left(\frac{\log\alpha}{\log\log\alpha}\right).
\]
Hence no such sequences exist.
\end{corollary}
\begin{proof}
Put $r=\rem(\beta,\alpha)$, so $r\in S_1$, $0\leq r<\alpha$ and
$\gcd(\alpha,r)=\gcd(\alpha,\beta)$. In a derivation from $\alpha,r$, the
value of the $k$th instruction lies in $S_{k+1}$, by induction on $k$: its
operands are $\alpha$, $r$ or earlier values, all in $S_k$. The source's
remainder agrees with ours at positive divisors. Theorem~\ref{thm:main}
bounds $G(\alpha,r)$.
\end{proof}

\section{Certificates and uniform processes}\label{sec:certificates}
\paragraph{Structures and certificates.}
Write $\mathrm{tt},\mathrm{ff}$ for the truth values and $\eq_c$ for the
test $x=c$. We use the book's partial remainder operation, defined only
when its second argument is nonzero, in the structures
\[
\mathcal N_0=(\mathbb N,\rem,\eq_0),\qquad
\mathcal N_\varepsilon=(\mathbb N,\rem,\eq_0,\eq_1).
\]
The latter is the book's Euclidean structure $\mathbf N_\varepsilon$
\cite[(1C-10), p.~24]{moschovakis2019}; our $\mathcal N_0$ should not
be confused with the book's $\mathbf N_0$, which has no remainder operation
\cite[p.~171]{moschovakis2019}.
Coprimeness is the relation
$x\coprime y\iff x,y\geq1\wedge\gcd(x,y)=1$
\cite[\S1A]{moschovakis2019}.

A derivation records finitely many remainder equations. These form a
partial structure generated by the inputs. A \emph{certificate} for $f$ at $\vec x$ in a structure
$\mathcal A$ is such a finite generated partial substructure, possibly with
test facts, on which every homomorphism into $\mathcal A$ respects $f$ at
$\vec x$. In other words, any assignment of values preserving the recorded
operations and tests must preserve the answer. The \emph{intrinsic
remainder-call complexity} $\calls_{\rem}(\mathcal A,f,\vec x)$ is the least
number of remainder facts in a certificate
\cite[\S\S1D, 4C--4E]{moschovakis2019}.

The Lean development represents a certificate by its facts. The following
lemma justifies this, and supplies the derivations used below.

\begin{lemma}[Fact representation]\label{lem:fact-representation}
Let $F$ be a finite set of true remainder and test facts in which every
number is generated from $a,b$ by the remainder facts of $F$. Then $F$ is
the diagram of a finite partial substructure generated by $a,b$ (of
$\mathcal N_\varepsilon$, or of $\mathcal N_0$ if $F$ has no $\eq_1$
facts), and every such substructure arises in this way. A map
$\pi:\mathbb N\to\mathbb N$ preserves the facts of $F$ exactly when its
restriction to the universe is a homomorphism, and every homomorphism
extends to such a map. Every element of the universe is the last value of
a derivation from $a,b$ with at most as many instructions as $F$ has
remainder facts.
\end{lemma}
\begin{proof}
Take the universe $V$ to consist of $a,b$ and the numbers occurring in $F$,
and interpret each primitive only at the recorded arguments, with the
recorded value. Truth of the facts makes these partial operations and
tests restrictions of the primitives of $\mathbb N$, and the hypothesis says that $V$ is generated by
$a,b$. Conversely, the diagram of a finite generated partial substructure is
such a set $F$. A homomorphism need only respect the recorded operations and
tests on $V$, and values outside $V$ are irrelevant; this gives the second
assertion. For the last, start from $a,b$ and repeatedly apply a remainder
fact of $F$ whose operands are available and whose value is new. While some
element of $V$ is unavailable such a fact exists: take the first
unavailable value in a generation of that element. Each fact is applied at
most once.
\end{proof}

\begin{proposition}\label{prop:certificates}
For coprime $a,b\geq1$,
$\calls_{\rem}(\mathcal N_\varepsilon,\coprime,a,b)=D(a,b)$. For
$(a,b)\neq(0,0)$,
\[
G(a,b)\leq\calls_{\rem}(\mathcal N_0,\gcd,a,b)\leq G(a,b)+2,
\]
and the same bounds hold for $\calls_{\rem}(\mathcal N_\varepsilon,\gcd,a,b)$.
\end{proposition}
\begin{proof}
A derivation of $1$, together with $\eq_1(1)=\mathrm{tt}$,
$\eq_0(a)=\mathrm{ff}$ and $\eq_0(b)=\mathrm{ff}$, certifies $\coprime$: under
any homomorphism the input images are positive, and each of their common
divisors divides every transported generated value, hence divides the image
$1$. Conversely, if a certificate has no positive $\eq_1$ fact, then
$x\mapsto2x$ preserves its remainder facts, zero tests and negative
one-tests, but not coprimeness. So every certificate contains $1$, and
Lemma~\ref{lem:fact-representation} gives a derivation of $1$ with at
most as many instructions as the certificate has remainder facts. This
proves the first assertion.

For gcd as a function, the certificate must reach its output $g$, giving a
lower bound of $G(a,b)$. For an upper bound, append
\[
\rem(a,g)=z,\qquad \rem(b,g)=z,\qquad \eq_0(z)=\mathrm{tt},
\]
where the generated $z$ is zero. Every homomorphism $\pi$ has $\pi(z)=0$ and
$\pi(g)>0$ (by definedness of either appended remainder). Thus $\pi(g)$
divides both input images. Conversely every common divisor of those images
divides $\pi(g)$ by induction over the derivation. This proves the bounds
for $\mathcal N_0$. The lower-bound argument applies to every gcd
certificate, $\eq_1$ facts included, and allowing $\eq_1$ facts can only lower the count, so the upper
bound transfers to $\calls_{\rem}(\mathcal N_\varepsilon,\gcd,a,b)$.
\end{proof}
\begin{remark}\label{rem:noncoprime}
At $(0,0)$ the output is an input, and the substructure with universe
$\{0\}$ and no facts is a certificate.
For gcd restricted to positive inputs, appending $\eq_0(a)=\mathrm{ff}$ and
$\eq_0(b)=\mathrm{ff}$ keeps input images in the domain, at no remainder cost. For a positive
noncoprime input, adding $\eq_1(g)=\mathrm{ff}$ to the gcd certificate certifies noncoprimeness, so the sublogarithmic remainder-call
bound holds for both truth values of coprimeness as well as for the gcd. The
overhead is a fixed additive constant and does not change the leading
coefficient, so Corollary~\ref{cor:six} holds for all three complexities of
Corollary~\ref{cor:refute} below on positive inputs.
\end{remark}

\paragraph{A single uniform process.}
Proposition~\ref{prop:certificates}, Remark~\ref{rem:noncoprime}, and
Theorem~\ref{thm:main} bound each intrinsic complexity by
$O(\log a/\log\log a)$, uniformly for $a\geq b\geq1$.
The book's Theorem~4F.1 states that, whenever every input has a certificate,
one uniform process attains the intrinsic complexity at every input
\cite[pp.~149--150]{moschovakis2019}. Hence:

\begin{corollary}\label{cor:uniform}
For each of gcd and coprimeness, there is a single uniform process of
$\mathcal N_\varepsilon$ computing it on $\mathbb N^2$ with
$O(\log a/\log\log a)$ remainder calls, uniformly for $a\geq b\geq1$.
The gcd process can also be taken over $\mathcal N_0$.
\end{corollary}
\begin{proof}
The certificate bounds give the cost assertion by Theorem~4F.1. Every input
has a certificate: symmetry covers reversed inputs, and a zero-test fact
suffices when either input is zero. Thus the same process is defined on
all of $\mathbb N^2$.
\end{proof}
The corollary rules out the possibility raised after Corollary~6C.7 of the
book that the Main Conjecture holds for all uniform processes of
$\mathcal N_\varepsilon$ \cite[p.~188]{moschovakis2019}.

\paragraph{The conjectures refuted.}
Moschovakis's 2012 notes state:
\begin{quote}
\textbf{Conjecture 3.4} \cite[p.~19]{moschovakis2012}. There is a rational number
$r>0$ such that for infinitely many $a\geq b\geq1$,
$\calls_{\rem}(\mathcal N_\varepsilon,\coprime,a,b)\geq r\log a$.
\end{quote}
The book formulates one precise version of the Main Conjecture
\cite[p.~151]{moschovakis2019}: there are a set $\mathcal D\subseteq\mathbb N^2$
and a $K>0$ such that
\[
(x,y)\in\mathcal D\ \Longrightarrow\
\calls_{\rem}(\mathrm{Eucl},x,y)\leq L(x,y)\leq K\calls_{\rem}(\mathcal A,\gcd,x,y),
\tag{4F-5}
\]
where $L(x,y)=2\log\max(x,y)$ for $(x,y)\neq(0,0)$, $L(0,0)=0$, and
$\calls_{\rem}(\mathrm{Eucl},x,y)$ counts
the remainder calls of the Euclidean algorithm $\mathrm{Eucl}$ (the book writes
$\varepsilon$) on $(x,y)$. The book states (4F-5) for the Euclidean algorithm
in $\mathbf N_\varepsilon$, so its $\mathcal A$ is $\mathcal N_\varepsilon$;
we treat both $\mathcal A=\mathcal N_\varepsilon$ and $\mathcal A=\mathcal N_0$.
The book adds that $\mathcal D$ then satisfies its conditions (4F-3) and
(4F-4); the latter makes the intrinsic complexity unbounded on $\mathcal D$,
so $\mathcal D$ is infinite (for finite $\mathcal D$ the statement is
trivial). We refute the weaker statement obtained by dropping the first
inequality and requiring only that $\mathcal D$ be infinite.

\begin{corollary}\label{cor:refute}
The intrinsic complexities
\[
\calls_{\rem}(\mathcal N_\varepsilon,\coprime,a,b),\qquad
\calls_{\rem}(\mathcal N_0,\gcd,a,b),\qquad
\calls_{\rem}(\mathcal N_\varepsilon,\gcd,a,b)
\]
are $O(\log a/\log\log a)$, uniformly for $a\geq b\geq1$. Consequently Conjecture~3.4 of
\cite{moschovakis2012} is false, and so is the version (4F-5) of the Main
Conjecture with $\mathcal A=\mathcal N_0$ or $\mathcal A=\mathcal N_\varepsilon$:
for every $r>0$ only finitely many pairs $a\geq b\geq1$ satisfy
$\calls_{\rem}(\mathcal N_\varepsilon,\coprime,a,b)\geq r\log a$, and for every
$K>0$ only finitely many pairs $(x,y)\in\mathbb N^2$ satisfy
$L(x,y)\leq K\calls_{\rem}(\mathcal A,\gcd,x,y)$.
\end{corollary}
The corollary follows from Theorem~\ref{thm:main} and
Proposition~\ref{prop:certificates}; for noncoprime pairs, the bound on the
coprimeness complexity comes from Remark~\ref{rem:noncoprime}. For (4F-5),
pairs with $x<y$ or a zero entry are handled as in the proof of
Corollary~\ref{cor:uniform}; the latter need no remainder facts.

\paragraph{The limit of the homomorphism method.}\label{par:method}
The Homomorphism Test proves lower bounds on intrinsic complexity by
showing that every generated diagram below the proposed cost admits a
homomorphism that fails to preserve the answer
\cite[Lemma~4E.3]{moschovakis2019}. Such bounds apply to every uniform
process computing the function, hence to every recursive program
\cite[Proposition~4E.1, Corollary~4E.2]{moschovakis2019}.
By Corollary~\ref{cor:refute}, no lower bound on the intrinsic complexities
listed there can grow faster than $\log a/\log\log a$, so a logarithmic lower bound
for algorithms must use restrictions beyond uniformity, such as finite program
syntax or determinism.

\section{An almost-all lower bound}\label{sec:loglog}
The difficult pairs of \cite[Theorem~6C.5]{moschovakis2019} give
$\Omega(\log\log a)$ on a sparse set. An elementary Diophantine argument gives
the same order for almost all coprime pairs. Write each generated value as
$ua+wb$. On the $B$-regular pairs defined below, small coefficients cannot give
a small positive value, and each remainder call at most squares the
coefficient bound, up to a factor $2B$. Reaching $1$ therefore takes about
$\log\log a$ calls, except on a sparse set of pairs.

\begin{proposition}\label{prop:loglog}
For integers $d\geq1$ and $N\geq3$, at most
$40(2+\ln N)N^{2-2^{-d}}$ coprime pairs $2\leq b<a\leq N$ satisfy
$D(a,b)\leq d$. Hence, for each fixed $\varepsilon>0$, all but $o(N^2)$ of
these pairs satisfy $D(a,b)>(1-\varepsilon)\log\log N$.
\end{proposition}

\begin{proof}
For $B\geq1$ call $(a,b)$ \emph{$B$-regular} if
$|qb-pa|\geq a/(Bq)$ whenever $q\in\mathbb N$, $q\geq1$,
$p\in\mathbb Z$ and $qb\neq pa$. We first bound the cost on regular
pairs, then count the irregular ones.

Every generated value has an integer representation $v=ua+wb$.
On a regular pair, $v>0$ implies $v\geq a/(B(|u|+|w|))$:
for $w\neq0$ use regularity with $q=|w|$, and for $w=0$ note $v\geq a$.
Let $H_t$ bound the coefficient norm $|u|+|w|$ in one such representation
of each value available after $t$ instructions, starting with $H_0=1$.
In an instruction $r=\rem(x,y)$, write $x=sy+r$. Since generated values
are at most $a$, the quotient satisfies $s\leq a/y\leq BH_t$.
Subtracting $s$ times the representation of $y$ from that of $x$ gives
\[
H_{t+1}\leq H_t+BH_t^2\leq2BH_t^2,\qquad
H_t\leq(2B)^{2^t-1}.
\]
Producing $1$ within $d$ instructions therefore forces
\[
1\geq\frac{a}{BH_d},\qquad a\leq\frac{(2B)^{2^d}}2.
\]

To count irregular pairs below $N$, choose a witness
$0<|qb-pa|<a/(Bq)$. Dividing $p,q$ by their gcd preserves this inequality,
so assume $\gcd(p,q)=1$. The nonzero integer $r=qb-pa$ then gives
$q<a/B\leq N/B$ and $0\leq p\leq q$.
For each $q$, there are at most $2q$ choices of $p$ and $2N/(Bq)$ choices
of the signed integer $r$. Given $p,q,r$, the equation
$pa\equiv-r\pmod q$ determines one residue class for $a$, leaving at most
$N/q+1$ choices; $b$ is then determined. With $Q=\lfloor N/B\rfloor$,
the number of irregular pairs is at most
\[
\begin{aligned}
\sum_{q=1}^{Q}(2q)\frac{2N}{Bq}\left(\frac Nq+1\right)
 &=\frac{4N^2}{B}\sum_{q=1}^{Q}\frac1q+\frac{4NQ}{B}\\
 &\leq\frac{4N^2}{B}(2+\ln N),
\end{aligned}
\]
using $Q\leq N$, $B\geq1$, and the harmonic-sum bound; an empty sum is zero.

Now split into dyadic ranges $X/2<a\leq X$, and take
$B=X^{2^{-d}}/2$. If $B\geq1$, a regular pair with $D\leq d$ would have
$a\leq X/2$, so all such pairs in the range are irregular. Their number is
at most $8(2+\ln X)X^{2-2^{-d}}$.
When $B<1$, the same bound follows from the trivial bound $X^2$, since
$X^{2^{-d}}<2$.
Put $s=2-2^{-d}\geq3/2$. Summing over the nonempty ranges with
$X=N/2^j\geq3$, and bounding their logarithmic factors by $2+\ln N$, gives
\[
8(2+\ln N)N^s\sum_{j\geq0}2^{-js}
 =\frac{8}{1-2^{-s}}(2+\ln N)N^s
 \leq40(2+\ln N)N^s.
\]
For the second assertion, take
$d=\lfloor\log\ln N-2\log\ln\ln N\rfloor$ for sufficiently large $N$.
Then $N^{-2^{-d}}\leq\exp(-(\ln\ln N)^2)$, so the exceptional set has
size $o(N^2)$, while $d=(1-o(1))\log\log N$, so eventually
$d\geq(1-\varepsilon)\log\log N$.
\end{proof}
A positive proportion (asymptotically $6/\pi^2$) of these pairs are coprime,
so the exceptional set also has density zero among coprime pairs.
The saving $2^{-d}$ in the exponent of $N^{2-2^{-d}}$ shrinks as $d$ grows,
so this argument does not reach the upper-bound scale $\log N/\log\log N$.

\section{Fixed programs with arbitrarily small coefficients}\label{sec:programs}

For a nondeterministic program $P$, let $c_P(a,b)$ be the minimum number of
remainder calls among its convergent computations. We require at least one
such computation on each input and require every convergent computation to
return the gcd. These are the conventions of
\cite[(2E-2), (2E-5)]{moschovakis2019}.

\begin{theorem}\label{thm:programs}
For each integer $p\geq1$, there is a finite nondeterministic recursive program
$P_p$ over $\mathcal N_0$ computing the gcd such that
\[
\limsup_{a\to\infty}\ \max_{1\leq b\leq a}
\frac{c_{P_p}(a,b)}{\log a}\leq\frac3p.
\]
Its generator, defined below, has $W=4^p+5$ individual parameters, and no
internal function has larger arity. Consequently, for every $\varepsilon>0$ there are a fixed
program $P$ and a constant $B\geq0$ such that
\[
\forall a\geq b\geq1,\qquad c_P(a,b)\leq\varepsilon\log a+B.
\]
\end{theorem}

The program depends on $\varepsilon$ but not on the inputs. The Main
Conjecture lets its coefficient $r$ depend on the program, so the family is
consistent with it; no single sublogarithmic program is obtained.

\paragraph{The program.}
Fix $W\geq3$ and write $\vec x=(x_0,\ldots,x_{W-1})$ and
$\vec x[i\leftarrow E]$ for the tuple with its $i$th expression replaced by
$E$. The generator can return a parameter, replace one by a remainder,
run a recursive subcomputation while retaining the caller's other values,
or copy parameters into new positions. These four actions are expressed
by the finite alternatives
\begin{align*}
\operatorname{choose}(\vec x)&=x_i &&(0\leq i<W),\tag{C}\\
\operatorname{gen}(\vec x)&=\operatorname{choose}(\vec x),\tag{G1}\\
\operatorname{gen}(\vec x)&=\operatorname{gen}
  (\vec x[i\leftarrow\rem(\operatorname{choose}(\vec x),
                               \operatorname{choose}(\vec x))])
  &&(0\leq i<W),\tag{G2}\\
\operatorname{gen}(\vec x)&=\operatorname{gen}
  (\vec x[i\leftarrow\operatorname{gen}(\vec x)]) &&(0\leq i<W),\tag{G3}\\
\operatorname{gen}(\vec x)&=\operatorname{gen}
  (\operatorname{choose}(\vec x),\ldots,\operatorname{choose}(\vec x)).\tag{G4}
\end{align*}
There are $W$ independent selector calls in (G4). Tuple notation and
indices are syntax abbreviations: after fixing $W$, each line expands into
ordinary first-order terms with fixed arities. No tuple, index, numeric
constant, or update operation is supplied as runtime data or a primitive.
The nested generator in (G3) is evaluated once; its caller's
other parameters remain available when it returns. The recursion stack can
grow, so $W$ is an arity bound, not a total storage bound.

Add $\operatorname{loop}(z)=\operatorname{loop}(z)$ and define
$\operatorname{accept}(a,b,g)$ by nested zero-test conditionals: if $g=0$,
loop; otherwise test $\rem(a,g)=0$ and then $\rem(b,g)=0$, returning $g$ if
both tests succeed and looping otherwise. The head is
\[
P_W(a,b)=\text{if }\eq_0(b)\text{ then }a\text{ else }
\operatorname{accept}(a,b,\operatorname{gen}(a,b,a,\ldots,a)).
\]
The four internal symbols have arities $W,W,3,1$ (for $\operatorname{choose}$,
$\operatorname{gen}$, $\operatorname{accept}$, $\operatorname{loop}$). Every common divisor of
the generator's parameters divides its output, by induction on a finite
evaluation. Thus a candidate is divisible by $\gcd(a,b)$; successful
verification proves the reverse divisibility. Every convergent computation
therefore returns the gcd. There is always a convergent computation following
Euclid, using parameter selection and copying to retain the current pair.
For $b>0$, final verification costs exactly two remainder calls.

This syntax is a finite recursive program in the sense of
\cite[\S\S2A, 2E]{moschovakis2019}; Appendix~\ref{app:lean} records the Lean
check that its evaluation trees correspond, with equal remainder cost, to
convergent computations of a specialisation of the book's recursive machine.

\begin{proof}[Proof of Theorem~\ref{thm:programs}]
We follow the construction of Section~\ref{sec:construction}: store the
prefix, contract, then recurse twice through nested generator calls.
Take $P_p=P_W$ with $W=4^p+5$; that is, set $k=2^p$, $h=k^2+1$ (the prefix
length, as in Corollary~\ref{cor:six} but now with $p$ fixed), and $W=h+4$.
Store $r_0,\ldots,r_h$ in $h+1$
parameters, reserving three parameters for the modulus and two scratch
residues. Lemmas~\ref{lem:gap} and~\ref{lem:prefix} give a contraction by a
factor greater than $k$ in at most three remainder calls, preserving the
prefix. The program can choose precisely the operands of each contraction.

The \emph{generator cost} of $(a,b)$ is the minimum number of remainder
calls in an evaluation of $\operatorname{gen}$ that returns $\gcd(a,b)$ from
$(a,b,a,\ldots,a)$. Let $T_W(n)$ be its maximum over $a,b\leq2^n$. Euclid
gives $T_W(n)=O(n)$. To \emph{route} to $x,y$ means to apply (G4) to pass to
the tuple $(x,y,x,\ldots,x)$, at no remainder cost. Copying parameters gives
symmetry. Multiplying every value by $g=\gcd(a,b)$ turns a run on
$(a/g,b/g)$ that returns $1$ into a run on $(a,b)$ that returns $g$, with the
same remainder count. It therefore suffices to analyse coprime inputs.

Fix $0<\beta<1/2$ and put $\alpha=1-\beta/2$. Use the early threshold
$N^\alpha$ and the target modulus $N^\beta$, where $N=2^n$. If no prefix
remainder is at most $N^\alpha$, then $r_h>N^\alpha$ and
$2a/r_h<2N^{1-\alpha}<N^\beta$ for all sufficiently large $n$.
There are at least $k+1$ distinct prefix residues modulo every modulus used
before stopping, since $h>k^2$. At most $\lceil n/p\rceil$ contractions
produce a positive $m\leq N^\beta$.

In the early case, route directly to the smaller consecutive Euclidean pair.
Otherwise form $u=a\bmod m$ and $v=b\bmod m$ in the scratch positions. A nested
generator call, with its parameters routed to copies of $m,u$, returns
$d=\gcd(m,u)$ while its caller retains $v$. Replace the modulus parameter
by $d$ and route to copies of $d,v$ for the second gcd computation. The
result is $\gcd(m,a,b)=1$. All these routes and nested calls are available
in the same fixed program. Thus, eventually in $n$,
\[
T_W(n)\leq\max\left\{
 T_W(\lceil\alpha n\rceil)+h+1,\quad
 2T_W(\lceil\beta n\rceil)+h+3\lceil n/p\rceil+2
\right\}.
\]
Set $\Lambda_W=\limsup_{n\to\infty}T_W(n)/n$, finite and nonnegative by Euclid.
Since $h$ is fixed, division by $n$ and passage to upper limits gives
\[
\Lambda_W\leq\max\{\alpha \Lambda_W,\ 2\beta \Lambda_W+3/p\}.
\]
As $\alpha<1$ and $\Lambda_W\geq0$, this implies $(1-2\beta)\Lambda_W\leq3/p$.
The program is independent of $\beta$, so letting $\beta\downarrow0$ gives $\Lambda_W\leq3/p$.
Take $n=\lceil\log a\rceil$ and add the two verification calls. Finally,
choose $p>3/\varepsilon$ and absorb the finitely many smaller inputs in $B$.
\end{proof}

\paragraph{Bounded memory.}
The nested calls in the proof are not needed. A \emph{register program}
keeps its values in a fixed number of registers and executes assignments
$x_l:=E$, branches on a test $E$, and a halt returning $E$, where each $E$ is
a term in the registers; the nondeterministic version allows several
commands at a line \cite[\S\S2E--2F]{moschovakis2019}. It is an iterator, so
it uses no recursion stack, and in pointed structures register programs
compute exactly the tail-recursive functions
\cite[Proposition~2F.1]{moschovakis2019}. Its cost $c_P$ counts remainder
evaluations as above.

\begin{theorem}\label{thm:registers}
For each integer $p\geq1$, there is a finite nondeterministic register
program over $\mathcal N_0$ with $4^p+7$ registers that computes the gcd on
inputs $a\geq b\geq1$ and satisfies the bound of
Theorem~\ref{thm:programs}:
\[
\limsup_{a\to\infty}\ \max_{1\leq b\leq a}\frac{c_P(a,b)}{\log a}\leq\frac3p.
\]
\end{theorem}

\begin{proof}
Put $k=2^p$, $h=k^2+1$ and $W=h+6$, and start the registers
$x_0,\ldots,x_{W-1}$ at $(a,b,a,b,a,\ldots,a)$. At its main line the program
may perform any assignment $x_l:=\rem(x_i,x_j)$ or $x_l:=x_i$ with $l\geq2$,
or pass to verification with candidate $x_i$. Verification halts with $x_i$
if $x_i\neq0$, $\rem(x_0,x_i)=0$ and $\rem(x_1,x_i)=0$, and loops otherwise.
The inputs $x_0,x_1$ are never overwritten. Every register value is
divisible by $\gcd(a,b)$, so every convergent computation returns the gcd,
and Euclid's algorithm gives one.

As before, scaling reduces the analysis to coprime pairs. Let $S(n)$ be the
largest, over coprime $1\leq b\leq a\leq2^n$, of the least number of
remainder calls that, starting from any state with $a,b$ in two working
registers, place $1$ in a working register. Fix $0<\beta<1/2$, put
$\alpha=1-\beta/2$ and $N=2^n$. Compute the Euclidean remainders
$r_0=a,r_1=b,r_2,\ldots$ in working registers, stopping at $r_h$ or at the
first $r_j\leq N^\alpha$. In the latter case one further remainder gives the
coprime pair $r_j,r_{j+1}$, bounded by $N^\alpha$, at total cost at most
$h$; continue with that pair. Otherwise $r_h>N^\alpha$, and as in the
proof of Theorem~\ref{thm:programs}, at most $\lceil n/p\rceil$
contractions, using a modulus register and two scratch registers, produce a
positive $m\leq N^\beta$. Then put $u=a\bmod m$ and $v=b\bmod m$, and run
Euclid's algorithm in place on $m,u$ and then on the result $d$ and $v$.
Each run takes at most $2\log m+3$ calls, and the result is
$\gcd(d,v)=\gcd(m,a,b)=1$. Thus, eventually in $n$,
\[
S(n)\leq\max\left\{S(\lceil\alpha n\rceil)+h,\quad
 h+3\lceil n/p\rceil+4\beta n+7\right\}.
\]
Euclid gives $S(n)=O(n)$. Passing to upper limits,
$\Lambda=\limsup_{n\to\infty}S(n)/n$ satisfies
$\Lambda\leq\max\{\alpha\Lambda,\ 3/p+4\beta\}$, and since $\alpha<1$,
$\Lambda\leq3/p+4\beta$. The program does not depend on $\beta$, so
$\Lambda\leq3/p$. Adding the two verification calls proves the bound.
\end{proof}
Here the completion uses Euclid's algorithm, whose cost $O(\beta n)$ is
negligible as $\beta\to0$; the recursive completion of
Section~\ref{sec:construction} is needed only when the prefix length grows.

\section{Discussion and open problems}\label{sec:discussion}
\paragraph{The gap.}
The intrinsic complexity of coprimeness in $\mathcal N_\varepsilon$ is at
least $\Omega(\log\log a)$ on almost all coprime pairs and at most
$O(\log a/\log\log a)$ on all positive pairs. Closing this gap is open;
unlike Bshouty's bound in his model, our upper bound is not known to be
optimal.

\paragraph{From a growing prefix to one program.}
The remaining obstacle in turning the construction into a sublogarithmic
program is reusable access to a prefix whose length grows with the input.
Selecting a close pair of residues costs nothing in a derivation, whereas
computing all prefix residues would take about $h$ calls per contraction.
The programs of Section~\ref{sec:programs} select operands nondeterministically
from a fixed number of registers. Can one fixed first-order recursive
program retain and repeatedly access a growing prefix at the same
remainder cost? Proposition~2E.1 of the book \cite[\S2E]{moschovakis2019}
converts each certificate into a computation of some program, not of one
program for all certificates.

\paragraph{Memory, not recursion.}
By Theorem~\ref{thm:registers}, $W$ registers give coefficient $3/p$ with
$W=4^p+7$, roughly $6/\log W$. So no bound $c\log a$ with $c$ independent
of the number of registers holds, even for nondeterministic register
programs. Does every nondeterministic register program with $W$ registers
make at least $c\log a/\log W$ remainder calls on infinitely many inputs,
for an absolute $c>0$? Such a bound cannot rest on how quickly values
shrink, since $\rem(m+1,m)=1$. Pratt's algorithm is itself a
nondeterministic register program \cite[Theorem~2E.2]{moschovakis2019};
compare also the book's question whether recursion reduces to tail
recursion in the pointed Euclidean structure
\cite[Problem~x2G.6]{moschovakis2019}.

\paragraph{Stronger lower-bound methods.}
The limit established in Section~\ref{sec:certificates} leaves room for
arguments combining homomorphisms with additional restrictions on
computations. In particular, the conjecture of a logarithmic lower bound
for deterministic uniform processes remains open
\cite[Conjecture~3.9, p.~23]{moschovakis2012},
\cite[\S4H and p.~188]{moschovakis2019}.

\section{Computational checks}\label{sec:checks}
The Python planner \nolinkurl{remainder_only_certificate.py}, in the
repository \url{https://github.com/jasondavies/euclid-uniform-processes},
chooses instructions offline with unrestricted arithmetic, using the
parameters of Section~\ref{sec:construction}; below a base size (256 bits by
default) it uses Euclid's algorithm. An independent replay then checks
operand availability, nonzero divisors, every remainder equation and the
final gcd.

The script \nolinkurl{scripts/benchmark_certificates.py} uses seed
$20261016$ to sample $25$ coprime pairs at each size: $a$ is uniform among
integers of that bit length, and $b$ is drawn uniformly from
$\{2,\ldots,a-1\}$ until $\gcd(a,b)=1$. It also runs the successive
Fibonacci pair $(F_k,F_{k-1})$ with $F_k$ the first Fibonacci number of
that bit length ($k=738$, $2951$, $11801$). The table gives remainder calls;
the Euclidean count stops when $1$ first appears.
\begin{center}
\small
\begin{tabular}{rcccc}\toprule
 & \multicolumn{2}{c}{Random pairs: min / median / max}
 & \multicolumn{2}{c}{Fibonacci pair}\\
\cmidrule(lr){2-3}\cmidrule(l){4-5}
Input bits & Construction & Euclid & Construction & Euclid\\\midrule
512 & 189 / 214 / 226 & 266 / 295 / 337 & 38 & 735\\
2048 & 617 / 647 / 669 & 1136 / 1192 / 1260 & 60 & 2948\\
8192 & 2007 / 2037 / 2059 & 4664 / 4774 / 4890 & 110 & 11798\\\bottomrule
\end{tabular}
\end{center}
All $78$ derivations and the example of Section~\ref{sec:contractions} pass
replay, and \nolinkurl{validation/benchmark_certificates.json} records the
individual counts. On the Fibonacci pairs, Euclid's worst case, the bits
removed by successive contractions roughly double (16, 31, 64, 125 and 252
at 512 bits), so few contractions are needed. These samples illustrate the construction; the
asymptotic claim rests on the proof. The
accompanying repository also checks scaled and noncoprime inputs, rejects invalid traces, and exhaustively tests the prefix
lemma over a finite range. For Section~\ref{sec:programs}, a separate compiler
and interpreter replay successful choices in finite first-order programs,
evaluating arguments right to left as the book's machine does. The
\texttt{README.md} records commands, test ranges, seeds, and results.
The Python interpreter is not formally verified; unrestricted arithmetic
used to plan choices is outside the programs.

\appendix
\section{Formal verification}\label{app:lean}
The development uses Lean~4.35.0 \cite{deMouraUllrich2021} with mathlib
\cite{mathlib2020} at its v4.35.0 tag; the pinned mathlib commit is
\texttt{56a005855ea17e0b0958f73ec801ae140b0a1679}.
The proved library contains no proof placeholders or added axioms; the main-result audit
reports only \texttt{propext}, \texttt{Classical.choice}, and
\texttt{Quot.sound}. Table~\ref{tab:lean} gives the main declarations;
the repository README gives the
complete map and reproduction commands.
The general theory of uniform processes, including Theorem~4F.1 used in
Corollary~\ref{cor:uniform}, and the historical lower bounds are cited,
not formalised. The certificate bounds and conjecture refutations are
formalised independently of that theory, with certificates represented by
their facts as in Lemma~\ref{lem:fact-representation}; that lemma,
Corollary~\ref{cor:closure-depth} and Theorem~\ref{thm:registers} are
proved here, not yet formalised.

The development is also packaged for the Palomar registry:
\texttt{Challenge.lean} imports only mathlib and states the main theorems
with every definition visible and deliberate \texttt{sorry} placeholders.
\texttt{Solution.lean} imports the proved library, independently of these placeholders.
Palomar's comparator checks that the proved statements are exactly those of
\texttt{Challenge.lean} and replays the proofs through Lean's kernel and
independent kernels. The thirteen compared theorems are marked ${}^\dagger$ in
Table~\ref{tab:lean}. The registry entry \cite{palomar2026} records this
verification for commit \texttt{40279dd}.

\begin{table}[htbp]
\centering\small
\caption{Main results and Lean declarations, all in the namespace \texttt{ShortGCD}. Names marked ${}^\dagger$ are compared against \texttt{Challenge.lean}.}
\label{tab:lean}
\begin{tabular}{@{}>{\raggedright\arraybackslash}p{0.29\textwidth}>{\raggedright\arraybackslash}p{0.67\textwidth}@{}}
\toprule
Paper & Lean declaration(s)\\\midrule
Theorem~\ref{thm:main} & \lean{uniform_real_bound}, \lean{sizeCost_isBigO}\\
Normalisation \eqref{eq:normalize} & \lean{G_normalize}\\
Lemma~\ref{lem:gap} & \lean{small_remainder}\\
Lemma~\ref{lem:prefix} & \lean{prefix_pairs_distinct}, \lean{residue_card_sq}\\
Corollary~\ref{cor:six} & \lean{uniform_real_bound_six}${}^\dagger$, \lean{uniform_intrinsic_bound_six}${}^\dagger$\\
Proposition~\ref{prop:certificates} &
\lean{callsRem_eq_D}, \lean{G_eq_D}, \lean{callsCoprimeness_eq_callsRem},
\lean{gcd_certificate_lower}, \lean{callsGcd_le}, \lean{callsRemGcd_le_callsGcd},
\lean{exists_certifiesGcd}${}^\dagger$, \lean{G_le_callsRemGcd}${}^\dagger$,
\lean{G_le_callsGcd}${}^\dagger$, \lean{callsCoprimeness_eq_G}${}^\dagger$\\
Remark~\ref{rem:noncoprime} & \lean{callsCoprimeness_le}, \lean{exists_certifiesCoprimeness}${}^\dagger$\\
Corollary~\ref{cor:uniform} & Cited Theorem~4F.1; not formalised here\\
Corollary~\ref{cor:refute} &
\lean{not_conjecture_3_4}${}^\dagger$, \lean{not_main_conjecture_4F5}${}^\dagger$,
\lean{not_main_conjecture_4F5_Neps}${}^\dagger$; literal versions
\lean{not_conjecture_3_4_literal}${}^\dagger$,
\lean{not_main_conjecture_4F5_literal}${}^\dagger$\\
Proposition~\ref{prop:loglog} & \lean{AlmostAll.card_D_le},
\lean{AlmostAll.tendsto_card_within_loglog},
\lean{tendsto_card_callsCoprimeness_loglog}${}^\dagger$\\
Theorem~\ref{thm:programs}: cost &
\lean{RegisterProgram.maxCost_step}, \lean{RegisterProgram.Fixed.limsup_profile_le}\\
Theorem~\ref{thm:programs}: syntax &
\lean{RegisterProgram.Syntax.eval_head_correct},
\lean{RegisterProgram.Syntax.eval_head_exists},
\lean{RegisterProgram.Syntax.arbitrarily_small_program}\\
Operational correspondence &
\lean{RegisterProgram.Syntax.Machine.eval_iff_steps},
\lean{RegisterProgram.Syntax.Machine.arbitrarily_small_program}\\
\bottomrule
\end{tabular}
\end{table}

\paragraph{Arithmetic bounds.}
Costs are lengths of explicit remainder derivations. The finite theorem
\lean{sizeCost_step} proves the recurrence from natural-number size
parameters; it does not assume a recurrence for the cost function.
For the initial Big-O proof, Lean uses
$j=\lfloor\log n\rfloor$, $p=\lfloor j/4\rfloor$ and $h=2^{2p}+1$
in place of $h=\lceil\sqrt n\rceil$, and checks the integer rounding
estimates before passing to real logarithms. The refinement uses the
tunable parameters of Corollary~\ref{cor:six} and checks the rounding limits
preserving coefficient $6$.
For Proposition~\ref{prop:loglog}, Lean takes $B=a^{2^{-d}}/2$ separately
for each pair instead of splitting into dyadic ranges. It proves the
stated constant $40$ and the exceptional-set limit for every fixed
$\varepsilon>0$, with threshold $(1-\varepsilon)\log\log N$.

\paragraph{Programs and machine computations.}
Generator runs translate into evaluations of finite first-order syntax
with exactly the same remainder count, including final verification.
Every terminating evaluation is correct, and a successful evaluation exists
on every input. Lean also proves exact-cost equivalence with a specialisation
of the book's two-stack machine: arguments are evaluated right to left,
zero tests and branches are separate transitions, and only remainder
transitions are charged \cite[Table~1, pp.~75--76]{moschovakis2019}.
The forward proof embeds each argument evaluation in its calling context;
the reverse proof decomposes a convergent computation into its argument
and body evaluations. Closures represent substitution in machine states,
not extra program data. This checks the language and program family used
here; the repository README describes the Python compiler and interpreter
and their relation to the Lean syntax.

\section*{Declaration of AI-generated research and authorship}
This manuscript is the product of an AI-driven research process. During a
broader investigation initiated by the human author, an AI system identified
the principal result. The author then instructed it to formalise the result in
Lean and to prepare this manuscript. The AI systems developed the result and
its presentation, including the mathematical exposition, the formalisation,
and the organisation and writing of the text; the human author wrote no part
of the manuscript. The research, the Lean development and the manuscript were
produced with GPT-6-Astra (OpenAI), used through Codex. Claude Opus~5.5
(Anthropic), used through Claude Code, carried out review and editing passes
and packaged the formalisation for the Palomar registry, which included
writing \texttt{Challenge.lean}, the lower-bound module and a proof change
needed by the independent kernels. GPT-6-Astra, through Codex, later revised
the exposition and expanded the computational checks.

The Lean development gives a machine-checkable version of the main claims,
subject to its trusted foundations (Lean's kernel, mathlib, and the axioms
listed in Appendix~\ref{app:lean}) and to the correspondence between its
definitions and the book's notions (Lemma~\ref{lem:fact-representation}).

This declaration distinguishes the provenance of the work from conventional
AI-assisted writing. The human author accepts responsibility for the
submission and for the accuracy of this disclosure.

\begin{thebibliography}{10}
\bibitem{bshouty1989}
Nader H. Bshouty.
\emph{Euclidean GCD algorithm is not optimal}.
Research Report 1989-376-38, Department of Computer Science, University of
Calgary, December 1989.
\href{https://doi.org/10.11575/PRISM/30475}{doi:10.11575/\allowbreak PRISM/30475}.

\bibitem{deMouraUllrich2021}
Leonardo de Moura and Sebastian Ullrich.
The Lean~4 theorem prover and programming language.
In \emph{Automated Deduction -- CADE~28}, Lecture Notes in Computer Science,
vol.~12699, pp.~625--635. Springer, 2021.
\href{https://doi.org/10.1007/978-3-030-79876-5_37}{doi:10.1007/978-3-030-79876-5\_37}.

\bibitem{knuth1997}
Donald E. Knuth.
\emph{The Art of Computer Programming, Vol.~2: Seminumerical Algorithms}, 3rd
ed. Addison-Wesley, Reading, MA, 1997.

\bibitem{mathlib2020}
The mathlib Community.
The Lean mathematical library.
In \emph{Proceedings of the 9th ACM SIGPLAN International Conference on
Certified Programs and Proofs (CPP 2020)}, pp.~367--381. ACM, 2020.
\href{https://doi.org/10.1145/3372885.3373824}{doi:10.1145/3372885.3373824}.

\bibitem{moschovakis2012}
Yiannis N. Moschovakis.
\emph{Absolute lower bounds for uniform processes}.
Unpublished lecture notes (first draft), January 25, 2012.
\url{https://www.math.ucla.edu/~ynm/nadalb.pdf}.

\bibitem{moschovakis2019}
Yiannis N. Moschovakis.
\emph{Abstract Recursion and Intrinsic Complexity}.
Lecture Notes in Logic, vol.~48. Cambridge University Press and Association
for Symbolic Logic, 2019.
\href{https://doi.org/10.1017/9781108234238}{doi:10.1017/9781108234238}.

\bibitem{palomar2026}
Palomar registry.
Entry PALOMAR-2026-10-23-000012, version~1: \emph{The Euclidean algorithm is
not optimal among uniform processes}. Registered October 23, 2026, for commit
\texttt{40279dd42196cc2f9dd6c14ecc3844b6822bff44} of the repository.
\url{https://palomar-registry.org/entry?id=PALOMAR-2026-10-23-000012&version=1}.

\bibitem{stein1967}
J. Stein.
Computational problems associated with Racah algebra.
\emph{Journal of Computational Physics} 1(3):397--405, 1967.
\href{https://doi.org/10.1016/0021-9991(67)90047-2}{doi:10.1016/0021-9991(67)90047-2}.

\bibitem{vddm2004}
Lou van den Dries and Yiannis N. Moschovakis.
Is the Euclidean algorithm optimal among its peers?
\emph{Bulletin of Symbolic Logic} 10(3):390--418, 2004.
\href{https://doi.org/10.2178/bsl/1102022663}{doi:10.2178/bsl/1102022663}.

\bibitem{vddm2009}
Lou van den Dries and Yiannis N. Moschovakis.
Arithmetic complexity.
\emph{ACM Transactions on Computational Logic} 10(1), Article~2, 49~pp., 2009.
\href{https://doi.org/10.1145/1459010.1459012}{doi:10.1145/1459010.1459012}.
Preprint: \url{https://www.math.ucla.edu/~ynm/papers/ac.pdf}.
\end{thebibliography}
\end{document}
